%=====================================================================
\section*{Problem 5. A prediction question}

\lead{Setup.} Take the looped d8w512 model of Problem 4: a prelude block, one core block applied $R=6$ times, and a coda block (15.6M parameters), trained for 9,375 updates on the A5 recipe (batch 64 sequences of 1,024 tokens, AdamW with $\beta=(0.9,0.95)$, weight decay 0.1, peak LR 0.003, 1\% warmup then linear decay, model seed 42, data seed 42, four micro-batches). The plain loop reaches 3.064 on 64 held-out sequences when run for the 6 passes it was trained with and 4.214 when run for 24 passes.

Now change two lines of the loop (diff below): apply a parameter-free RMSNorm to the residual stream at every pass boundary and add the prelude output back in afterwards (``norm + recall''). Train it identically and evaluate it at \textbf{24 passes}, four times the trained count.

\lead{Question.} At 24 passes the held-out loss of the norm+recall loop is
\begin{enumerate}[label=(\Alph*)]
\item above 4.5 (worse than the plain loop: the norm removes scale information and recall keeps re-injecting a stale signal),
\item between 3.5 and 4.5 (about the same degradation as the plain loop),
\item between 3.1 and 3.5 (degrades, but much less than the plain loop),
\item below 3.1 (within 0.1 of its loss at the trained 6 passes).
\end{enumerate}

\lead{Answer: (D).} Measured: 3.033 at 24 passes against 3.023 at the trained 6 passes, and 3.035, 3.030, 3.032, 3.033 at 8, 9, 12 and 16 passes. The matched control is the plain six-pass loop trained with the same recipe and seed: 3.064 at 6 passes, 3.113 at 8, 3.349 at 12, 3.794 at 18 and 4.214 at 24. This agrees with our prediction. The norm turns each pass into a map on a sphere of fixed radius, and recall gives it a fixed point to settle on; after about six applications the stream stops moving, so additional passes are free. Without recall the same norm collapses training to the unigram loss (Problem 4(b)), and recall without the norm degrades like the plain loop (3.40 at 16 passes), so both lines are needed.

\lead{Diff} against the plain \texttt{LoopLM} in \texttt{p4\_looping.py} (14 lines):
\begin{verbatim}
 class LoopLM(AutoregressiveLM):
-    def __init__(self, config, *, passes=6, core_size=1, **kwargs):
+    def __init__(self, config, *, passes=6, core_size=1,
+                 boundary_norm=False, recall=0., **kwargs):
         super().__init__(config, **kwargs)
         self.passes, self.core_size = passes, core_size
+        self.boundary = (nn.RMSNorm(config.hidden_size, eps=1e-5, elementwise_affine=False)
+                         if boundary_norm else nn.Identity())
+        self.recall = recall

     def forward_hidden(self, h, rope, attention_mask):
         h = self.model.layers[0](h, rope, attention_mask)   # prelude
+        anchor = h
         for r in range(self.passes):
             for j in range(self.core_size):
                 h = self.model.layers[1 + j](h, rope, attention_mask)
+            h = self.boundary(h) + self.recall * anchor
         return self.model.layers[-1](h, rope, attention_mask)  # coda
\end{verbatim}
Reproduce with \texttt{config('loop-norm-recall')} (\texttt{passes=6, core\_size=1, boundary\_norm=True, recall=1.0}) and the control \texttt{config('loop-r6')} in \texttt{sweeps.py}; \texttt{measure\_runs.py} evaluates a saved checkpoint at 1 to 24 passes on the first 64 validation sequences.
